% chktex-file 3
% chktex-file 35
\section{Causal protection and the observer interior}\label{sec:causal-protection}

\subsection{Exact comoving plateau geometry}

Let the CausalPoint observer be placed in an open interior region
\(\mathcal U_{\mathrm{int}}\) on which the SHBT shape functional has an exact
plateau.  The defining conditions are
\begin{equation}
  f_{\mathrm{SHBT}}(\boldsymbol{x},\theta)
  =
  1,
  \qquad
  \partial_i f_{\mathrm{SHBT}}
  =
  0,
  \qquad
  \beta_x(\boldsymbol{x})
  =
  \beta_0
  =
  -c
  \exp\left(
    \frac{\Delta_{\mathrm{mod}}}{2}
  \right).
  \label{eq:exact-interior-plateau}
\end{equation}
Because the plateau is open and constant, all higher spatial derivatives of
\(\beta_x\) vanish in \(\mathcal U_{\mathrm{int}}\).  The four-dimensional
line element induced by Eq.~\eqref{eq:four-dimensional-shift-line-element}
therefore reduces to
\begin{equation}
  ds_{\mathrm{int}}^2
  =
  -c^2dt^2
  +
  \left(
    dx+\beta_0dt
  \right)^2
  +
  dy^2
  +
  dw^2.
  \label{eq:interior-plateau-metric}
\end{equation}
Introduce the comoving coordinates
\begin{equation}
  T=t,
  \qquad
  X=x+\beta_0t,
  \qquad
  Y=y,
  \qquad
  W=w.
  \label{eq:interior-coordinate-transformation}
\end{equation}
Since \(\beta_0\) is constant,
\begin{equation}
  dT=dt,
  \qquad
  dX=dx+\beta_0dt,
  \qquad
  dY=dy,
  \qquad
  dW=dw.
  \label{eq:interior-coordinate-differentials}
\end{equation}
Substitution into Eq.~\eqref{eq:interior-plateau-metric} gives
\begin{equation}
  ds_{\mathrm{int}}^2
  =
  -c^2dT^2
  +
  dX^2
  +
  dY^2
  +
  dW^2
  =
  \eta_{\hat\mu\hat\nu}
  dX^{\hat\mu}dX^{\hat\nu}.
  \label{eq:exact-interior-minkowski}
\end{equation}
Thus the plateau shift is a constant-coordinate shear, not an intrinsic
curvature of the interior.  Its magnitude may exceed \(c\) relative to the
external coordinates without changing the local Minkowski signature.

\subsection{Vanishing connection, curvature, and proper acceleration}

In the coordinates \(X^{\hat\mu}=(cT,X,Y,W)\), the metric components are
constant:
\begin{equation}
  g_{\hat\mu\hat\nu}^{\mathrm{obs}}
  =
  \eta_{\hat\mu\hat\nu}
  =
  \operatorname{diag}(-1,1,1,1),
  \qquad
  \partial_{\hat\alpha}
  g_{\hat\mu\hat\nu}^{\mathrm{obs}}
  =
  0.
  \label{eq:observer-metric-constant}
\end{equation}
The Levi-Civita connection is consequently
\begin{equation}
  \Gamma^{\hat\alpha}{}_{\hat\mu\hat\nu}
  =
  \frac{1}{2}
  g_{\mathrm{obs}}^{\hat\alpha\hat\sigma}
  \left(
    \partial_{\hat\mu}
    g_{\hat\sigma\hat\nu}^{\mathrm{obs}}
    +
    \partial_{\hat\nu}
    g_{\hat\sigma\hat\mu}^{\mathrm{obs}}
    -
    \partial_{\hat\sigma}
    g_{\hat\mu\hat\nu}^{\mathrm{obs}}
  \right)
  =
  0.
  \label{eq:interior-christoffel-zero}
\end{equation}
It follows directly that
\begin{equation}
  R^{\hat\alpha}{}_{\hat\beta\hat\mu\hat\nu}
  =
  \partial_{\hat\mu}
  \Gamma^{\hat\alpha}{}_{\hat\nu\hat\beta}
  -
  \partial_{\hat\nu}
  \Gamma^{\hat\alpha}{}_{\hat\mu\hat\beta}
  +
  \Gamma^{\hat\alpha}{}_{\hat\mu\hat\sigma}
  \Gamma^{\hat\sigma}{}_{\hat\nu\hat\beta}
  -
  \Gamma^{\hat\alpha}{}_{\hat\nu\hat\sigma}
  \Gamma^{\hat\sigma}{}_{\hat\mu\hat\beta}
  =
  0.
  \label{eq:interior-riemann-zero}
\end{equation}
All contractions of the Riemann tensor vanish as well:
\begin{equation}
  R_{\hat\mu\hat\nu}
  =
  0,
  \qquad
  R
  =
  0,
  \qquad
  G_{\hat\mu\hat\nu}
  =
  0
  \quad
  \text{on }
  \mathcal U_{\mathrm{int}}.
  \label{eq:interior-curvature-contractions}
\end{equation}

A comoving observer has fixed \((X,Y,W)\) and four-velocity
\begin{equation}
  u^{\hat\mu}
  =
  \frac{dX^{\hat\mu}}{d\tau_{\mathrm p}}
  =
  (c,0,0,0),
  \qquad
  \eta_{\hat\mu\hat\nu}
  u^{\hat\mu}u^{\hat\nu}
  =
  -c^2.
  \label{eq:comoving-interior-four-velocity}
\end{equation}
Its covariant four-acceleration is
\begin{equation}
  a^{\hat\mu}
  =
  u^{\hat\nu}
  \nabla_{\hat\nu}
  u^{\hat\mu}
  =
  u^{\hat\nu}
  \partial_{\hat\nu}
  u^{\hat\mu}
  +
  \Gamma^{\hat\mu}{}_{\hat\alpha\hat\beta}
  u^{\hat\alpha}u^{\hat\beta}
  =
  0.
  \label{eq:interior-proper-acceleration-zero}
\end{equation}
The executable observer audit returns
\begin{equation}
  \left\|
    g_{\hat\mu\hat\nu}^{\mathrm{obs}}
    -
    \eta_{\hat\mu\hat\nu}
  \right\|_{\infty}
  =
  \SimOutputObserverMetricError,
  \qquad
  \|a^{\hat\mu}\|
  =
  \SimOutputAccelerationNorm
  \;\mathrm{m\,s^{-2}}.
  \label{eq:numerical-observer-audit}
\end{equation}
The zero-acceleration result applies to the exact plateau.  The bubble wall,
where derivatives of \(f_{\mathrm{SHBT}}\) are nonzero, is not locally
Minkowskian and must be analyzed separately.

\subsection{Chronology from a global resolution time}

Monotonicity of a coordinate along one trajectory is not sufficient by itself
to exclude closed timelike curves.  The required global statement is that the
boundary resolution variable \(\tau\) is a smooth time function on the
admissible spacetime domain.  We impose
\begin{equation}
  \frac{d\tau}{dt}
  =
  H(t)
  >0,
  \qquad
  g^{\mu\nu}
  \nabla_\mu\tau
  \nabla_\nu\tau
  <0.
  \label{eq:global-resolution-time-condition}
\end{equation}
The second inequality states that \(\nabla_\mu\tau\) is everywhere timelike.
Choose its orientation so that \(\tau\) increases on every future-directed
causal curve \(x^\mu(\lambda)\).  Then
\begin{equation}
  \frac{d\tau}{d\lambda}
  =
  \frac{dx^\mu}{d\lambda}
  \nabla_\mu\tau
  >0.
  \label{eq:resolution-increases-on-causal-curves}
\end{equation}
Suppose, for contradiction, that a future-directed closed causal curve
\(\mathcal C\) exists with
\(x^\mu(\lambda_1)=x^\mu(\lambda_0)\).  Single-valuedness of the time function
would require
\begin{equation}
  \tau(\lambda_1)-\tau(\lambda_0)
  =
  0.
  \label{eq:closed-curve-time-return}
\end{equation}
Monotonicity instead gives
\begin{equation}
  \tau(\lambda_1)-\tau(\lambda_0)
  =
  \int_{\lambda_0}^{\lambda_1}
  \frac{d\tau}{d\lambda}
  d\lambda
  >0,
  \label{eq:closed-curve-contradiction}
\end{equation}
contradicting Eq.~\eqref{eq:closed-curve-time-return}.  Therefore
\begin{equation}
  \nexists\,
  \text{closed future-directed causal curve}
  \quad
  \text{whenever Eq.~\eqref{eq:global-resolution-time-condition} holds}.
  \label{eq:ctc-exclusion-theorem}
\end{equation}
This is the precise chronology-protection statement: the monotonic SHBT
resolution arrow suppresses CTCs provided it extends to a globally defined
time function with timelike gradient.